{\footnotesize
\setlength{\tabcolsep}{3.5pt}
\begin{threeparttable}
\begin{tabular}{lllrrrrr}
\toprule
 & Region & State & Initial employment & $T$ & $D$ & $F$ & Portab. \\
\midrule
Most flexible & Campinas & São Paulo & 1,837,228 & 0.042 & 0.016 & 2.65 & 0.698 \\
 & Porto Alegre & Rio Grande do Sul & 1,134,515 & 0.046 & 0.019 & 2.35 & 0.655 \\
 & Criciúma & Santa Catarina & 242,473 & 0.046 & 0.020 & 2.32 & 0.675 \\
 & Sorocaba & São Paulo & 624,235 & 0.038 & 0.017 & 2.20 & 0.663 \\
 & Chapecó & Santa Catarina & 305,789 & 0.037 & 0.017 & 2.18 & 0.672 \\
 & Rio de Janeiro & Rio de Janeiro & 2,417,403 & 0.054 & 0.025 & 2.15 & 0.664 \\
 & São Paulo & São Paulo & 6,522,499 & 0.047 & 0.022 & 2.12 & 0.696 \\
 & Ijuí & Rio Grande do Sul & 134,022 & 0.052 & 0.026 & 1.98 & 0.659 \\
 & Caçador & Santa Catarina & 61,932 & 0.052 & 0.026 & 1.98 & 0.667 \\
 & Caxias do Sul & Rio Grande do Sul & 364,042 & 0.036 & 0.019 & 1.93 & 0.690 \\
\addlinespace[3pt]
Least flexible & Lábrea & Amazonas & 4,444 & 0.147 & 1.661 & 0.09 & 0.575 \\
 & Gurupi & Tocantins & 23,021 & 0.103 & 1.199 & 0.09 & 0.583 \\
 & Crateús & Ceará & 9,117 & 0.161 & 2.641 & 0.06 & 0.584 \\
 & Picos & Piauí & 17,185 & 0.130 & 2.150 & 0.06 & 0.637 \\
 & Rio Branco & Acre & 62,778 & 0.106 & 1.865 & 0.06 & 0.643 \\
 & Paulo Afonso & Bahia & 24,263 & 0.125 & 2.355 & 0.05 & 0.595 \\
 & Corrente - Bom Jesus & Piauí & 7,106 & 0.123 & 3.885 & 0.03 & 0.607 \\
 & Cruzeiro do Sul & Acre & 8,205 & 0.094 & 3.242 & 0.03 & 0.590 \\
 & Imperatriz & Maranhão & 100,051 & 0.061 & 2.534 & 0.02 & 0.655 \\
 & Corumbá & Mato Grosso do Sul & 37,383 & 0.086 & 4.645 & 0.02 & 0.545 \\
\bottomrule
\end{tabular}
\begin{tablenotes}[flushleft]\footnotesize
\item \textit{Note:} Intermediate regions ranked by allocative flexibility in the 2019 to 2022 window. Source: RAIS, authors' calculations.
\end{tablenotes}
\end{threeparttable}
}
